\scanpage{159}
% archive foliation: 158

\uline{Trialité des bi-di-triades}

On peut voir la structure \unsure{des deux parties}, soit par un \unsure{ens.}: 9 éléments $V$ (plan affine sur $\mathbb F_3$), soit par l'autre: 12 éléments $E$ (ens des droites de $V$). Relations d'incidence entre $E$ et $V$.

\margin{NB $E$ est \unsure{un} plan projectif sur $\mathbb F_3$, privé d'un de ses points.}

A) La structure sur $V$ revient à la donnée d'une partie $E$ de $\mathrm P_3(V)$, telle que
\[
\left\{\begin{array}{l}
\operatorname{card}E=12\\
\forall a,b\in E,\ \text{on a}\quad \operatorname{card}a\cap b\leq1
\end{array}\right.
\]
\deleted{\unsure{(en fait}}

\deleted{\unsure{B)}}

B) Structure sur $E$: 1) partition en quatre classes de 3 éléments (\unsure{suivant} les quatre directions de droites)
\[
E=\coprod_{i\in I}E_i
\]

2) \deleted{Ens} $V$ de parties à 4 élts de $E$. (ens des droites passant par un point \inserted{de $V$}) avec les propriétés

a) Tout \deleted{\unsure{droite}} $D\in V$ rencontre chaque $E_i$ en \inserted{(au moins)} un \uline{pt} (\unsure{NB}: \inserted{au plus} un, \deleted{seul} \inserted{\unsure{revient au $=$}} \deleted{(de $D$)}
\[
D\longrightarrow I\quad\text{est bijectif})
\]

b) Si $D,D'\in V$, $D\neq D'$, alors $\operatorname{card}D\cap D'=1$

\margin{\uline{NB} b) résulte de a) et c)}

c) Si $a,b\in E$ n'appartiennent pas à un même $E_i$, $\exists$ unique $D\in V$ qui y passe. (unicité provient d'ailleurs de b))

\margin{En introduisant $\widehat E\inserted{=E\coprod\{\infty\}}$ et
\[
\widehat V=V\coprod\left\{\begin{array}{l}
\text{ens des parties de }\widehat E\text{ de la forme}\\
E_i\cup\{\infty\}
\end{array}\right\}
\]
Ceci peut se dire: \deleted{\unsure{chaque}} deux ``droites'' distinctes $(\in\widehat V)$ se rencontrent en un point unique.}

\uline{NB} c implique déjà que $\operatorname{card}V=\unsure{9}$, et compte tenu de a), que par $\forall$ point $a\in E$ passent \unsure{exact$^t$} 3 ``droites''.
